\documentclass[11pt,a4paper]{article}
\usepackage[utf8]{inputenc}
\usepackage{amsmath,amssymb,amsfonts}
\usepackage{geometry}
\geometry{margin=1in}

\title{The HQR Lagrangian and the Informational Stress-Energy Tensor \\ (Path A Formulation)}
\author{Holonomic Quantum Research Group}
\date{\today}

\begin{document}

\maketitle

\begin{abstract}
We present the 11-dimensional Lagrangian of Holonomic Quantum Reality under the
Path A commitment, in which \(\rho^{\rm ent}\) is treated as a fundamental
phenomenological scalar field. We fix an explicit metric signature, derive the
quantum informational stress-energy tensor \(T^{\rm info}_{AB}\) with internally
consistent kinetic and potential signs, and demonstrate explicit dimensional
reduction on a flat 7-torus. We also state precisely why the torus reduction is
a consistency check rather than a distinctive prediction.
\end{abstract}

\section{Path A Commitment}

In this formulation, \(\rho^{\rm ent}\) is a fundamental scalar field on the 11D
manifold. The name ``entanglement density'' is interpretive and does not imply
that \(\rho^{\rm ent}\) is derived from the von Neumann entropy of a reduced
density matrix.

\section{Conventions}

\textbf{Metric signature.} We use the mostly-plus convention
\((-,+,+,\dots,+)\) in \(D=11\). With this signature, a healthy (non-ghost)
scalar enters the Lagrangian with a \emph{negative} kinetic term and a
\emph{negative} potential:
\[
\mathcal{L}_{\rm scalar}
= -\tfrac{\beta}{2}\, g^{AB}\nabla_A\rho\,\nabla_B\rho \;-\; V(\rho),
\qquad \beta>0,
\]
so that, for a slowly varying field, \(\mathcal{L}_{\rm scalar}\to -V(\rho)\)
and the energy density is \(+V(\rho)\) (bounded below for \(V\ge 0\)). Fixing
this convention here is essential: the relative sign between the
gradient-squared and potential pieces of \(T^{\rm info}_{AB}\) is what
determines whether the informational sector behaves as dark energy or as
ordinary matter after reduction, and an inconsistent choice flips the physical
prediction.

\section{The 11D Lagrangian}

The Lagrangian density is
\begin{equation}
\mathcal{L}_{\rm HQR} = \frac{R^{(11)}}{16\pi G_{11}}
+ \alpha \, \rho \, R^{(11)}
- \frac{\beta}{2} g^{AB} \nabla_A \rho \, \nabla_B \rho
- \gamma \, V(\rho)
+ \mathcal{L}_{\rm hidden},
\end{equation}
where \(\rho = \rho^{\rm ent}\), \(\beta>0\), and \(\gamma>0\). Note the signs:
the kinetic and potential terms enter with the \emph{same} (negative) sign
demanded by the mostly-plus convention of Section~2. (This corrects an earlier
draft in which both appeared with a leading \(+\), which is either a ghost or a
sign error.)

The hidden-order sector includes topological contributions:
\begin{equation}
\mathcal{L}_{\rm hidden} = \frac{\kappa}{3!} C \wedge dC \wedge dC
+ \lambda \, \rho \, \operatorname{Tr}(R \wedge R \wedge R \wedge R \wedge R)
+ \theta(\rho) \, \operatorname{CS}_5(F).
\end{equation}
These are topological/boundary terms: they are metric-independent (or
total-derivative) and therefore do not contribute to the local metric variation
that defines \(T^{\rm info}_{AB}\); they instead affect the \(\rho\) equation of
motion and global selection rules. This is stated explicitly so that no
spurious local stress-energy is attributed to them downstream.

\section{Derivation of \( T^{\rm info}_{AB} \)}

Define \( f = \dfrac{1}{16\pi G_{11}} + \alpha \rho \). Varying the action with
respect to the metric (mostly-plus) yields the modified Einstein equation
\begin{equation}
f \big( R_{AB} - \tfrac12 R g_{AB} \big)
- \big( \nabla_A \nabla_B - g_{AB} \square \big)(\alpha \rho)
= 8\pi G_{11} \, T_{AB}^{\rm (kin+pot+matter)} .
\end{equation}
Collecting terms into standard form gives the effective informational
stress-energy tensor
\begin{equation}
T_{AB}^{\rm info} = \frac{1}{f} \Bigg[
\beta \Big( \nabla_A \rho \, \nabla_B \rho - \tfrac12 g_{AB} |\nabla\rho|^2 \Big)
\;-\; g_{AB}\, \gamma\, V(\rho)
\;+\; \frac{\alpha}{8\pi G_{11}} \big( \nabla_A \nabla_B \rho - g_{AB} \square \rho \big)
\Bigg].
\end{equation}

\paragraph{Sign check (the load-bearing point).}
With the mostly-plus convention of Section~2, the variation of
\(-\tfrac{\beta}{2}(\nabla\rho)^2\) produces the
\(\beta(\nabla_A\rho\nabla_B\rho - \tfrac12 g_{AB}|\nabla\rho|^2)\) structure
above, and the variation of \(-\gamma V(\rho)\) produces
\(-g_{AB}\,\gamma V(\rho)\). The gradient-squared and potential pieces therefore
carry \emph{opposite} relative signs (the gradient piece \(+\), the potential
piece \(-g_{AB}\)). This differs from the earlier draft, which wrote the
potential as \(+g_{AB}\gamma V(\rho)\) --- inconsistent with a Lagrangian in
which the kinetic and potential signs are not independently chosen. The
corrected form is required for the \(00\)-component to give a positive energy
density \(\propto +\gamma V(\rho)\) for \(V\ge 0\), and it is this sign that
later fixes whether the informational sector accelerates or decelerates the
effective 4D expansion. Any downstream cosmological or condensed-matter estimate
must use this corrected sign.

\section{Explicit Reduction on Flat \( T^7 \)}

Consider the product metric ansatz
\[
ds_{11}^2 = g_{\mu\nu}(x)\, dx^\mu dx^\nu + \delta_{mn}\, dy^m dy^n
\]
on \( M^4 \times T^7 \), with \(\rho = \rho_0(x)\) independent of the internal
coordinates.

The internal derivatives vanish, and the reduction of the non-minimal
contribution to the effective 4D energy density is
\[
T_{00}^{(4)}(x) = C \big( \nabla_0 \nabla_0 \rho_0 - g_{00} \square_4 \rho_0 \big),
\]
where \( C \) is a constant set by \(\alpha\) and \( G_{11}\). The volume
factors cancel exactly, yielding a finite, definite 4D expression.

\paragraph{What this does and does not establish.}
This demonstrates that the reduction \emph{procedure} is well-defined and
calculable once the compactification ansatz and the profile of \(\rho\) are
specified --- i.e. it is a consistency check on the machinery. It does
\textbf{not} constitute a distinctive HQR prediction. On a flat \(T^7\) with a
zero-mode profile, the surviving 4D structure is degenerate with that of an
ordinary non-minimally coupled modulus (a generic scalar--tensor / Brans--Dicke
theory). Nothing in this reduced expression distinguishes HQR from that generic
case. A distinctive prediction requires a \emph{structured} internal manifold
(warped product, non-trivial flux, or a \(G_2\)/Calabi--Yau geometry) in which
the topological hidden-order terms and the internal curvature feed back into the
4D effective theory. The torus is therefore a stress test of consistency only,
and any prediction extracted at this level (including for URu\(_2\)Si\(_2\))
would not be a test of HQR specifically.

\section{Notes}

\begin{itemize}
    \item The torus example is a minimal consistency check, not a physical
          prediction (see the paragraph above). More realistic compactifications
          (warped products, fluxes, \(G_2\) or Calabi--Yau manifolds) follow the
          same structural procedure but involve additional technical steps and
          are the only route to a distinctive signature.
    \item Topological terms in \(\mathcal{L}_{\rm hidden}\) are
          metric-independent/total-derivative and primarily affect the dynamics
          of \(\rho\) and global selection rules; they do not contribute
          directly to the local stress-energy tensor.
    \item \textbf{Open:} the dynamics/stabilisation of \(\rho\) (its equation of
          motion plus the topological selection rules) still has to fix which
          profile \(\rho_0(x)\) is realised. Until that is done, \(\rho_0\)
          remains free and no numerical prediction is fixed.
\end{itemize}

\end{document}
